\documentclass{article}

% Language setting
\usepackage[english]{babel}

% Set page size and margins
\usepackage[letterpaper,top=2cm,bottom=2cm,left=3cm,right=3cm,marginparwidth=1.75cm]{geometry}

% Useful packages
\usepackage{amsmath}
\usepackage{amssymb}
\usepackage{graphicx}
\usepackage[colorlinks=true, allcolors=black]{hyperref}
\usepackage{titlesec}
\usepackage{xparse}

\NewDocumentCommand{\qcq}{m o o m}{%
	\ensuremath{\forall #1 \IfValueT{#2}{, #2} \IfValueT{#3}{, #3}  \in  \mathbb{#4}}%
}
\NewDocumentCommand{\ext}{m o o m}{%
	\ensuremath{\exists #1 \IfValueT{#2}{, #2} \IfValueT{#3}{, #3}  \in  \mathbb{#4}}%
}

\title{\textbf{Structures Algébriques}}
\author{\textbf{Yassin Akkab}}

\newcommand{\R}{\mathbb{R}}
\newcommand{\N}{\mathbb{N}}
\newcommand{\Z}{\mathbb{Z}}
\newcommand{\C}{\mathbb{C}}
\newcommand{\Q}{\mathbb{Q}}
\newcommand{\D}{\mathbb{D}}

\begin{document}
	
	\maketitle
	\tableofcontents
	\newpage
	
	\section{Lois de composition interne (LCI)}
	
	\subsection{Définition}
	Soit $E$ un ensemble. Une \textbf{loi de composition interne} (LCI) sur $E$ est une application :
	\[ *: E \times E \to E, \quad (x, y) \mapsto x * y \]
	
	\subsection{Propriétés d'une LCI}
	Soit $*$ une LCI sur $E$.
	\begin{itemize}
		\item \textbf{Associativité} : La loi $*$ est associative si :
		\[ \qcq{x}[, y][, z]{E}, \quad (x * y) * z = x * (y * z) \]
		\item \textbf{Commutativité} : La loi $*$ est commutative si :
		\[ \qcq{x}[, y]{E}, \quad x * y = y * x \]
		\item \textbf{Élément neutre} : Un élément $e \in E$ est un élément neutre pour $*$ si :
		\[ \qcq{x}{E}, \quad x * e = e * x = x \]
		S'il existe, l'élément neutre est unique.
		\item \textbf{Élément symétrisable} : Si $E$ admet un élément neutre $e$, un élément $x \in E$ est dit symétrisable pour $*$ s'il existe $x' \in E$ tel que :
		\[ x * x' = x' * x = e \]
		$x'$ est appelé le symétrique de $x$. Si la loi est associative, le symétrique est unique.
	\end{itemize}
	
	\section{Groupes et sous-groupes}
	
	\subsection{Définition d'un groupe}
	Un \textbf{groupe} est un couple $(G, *)$ où $G$ est un ensemble et $*$ est une LCI sur $G$ vérifiant :
	\begin{enumerate}
		\item La loi $*$ est associative.
		\item La loi $*$ admet un élément neutre $e \in G$.
		\item Tout élément de $G$ admet un symétrique dans $G$ pour la loi $*$.
	\end{enumerate}
	Si de plus la loi $*$ est commutative, le groupe est dit \textbf{commutatif} (ou \textbf{abélien}).
	
	\subsection{Sous-groupes}
	Soit $(G, *)$ un groupe d'élément neutre $e$. Une partie non vide $H$ de $G$ est un \textbf{sous-groupe} de $(G, *)$ si $(H, *)$ est lui-même un groupe.
	\textbf{Caractérisation d'un sous-groupe} :
	$H$ est un sous-groupe de $(G, *)$ si et seulement si :
	\begin{enumerate}
		\item $H \ne \emptyset$ (en particulier, $e \in H$).
		\item $\qcq{x}[, y]{H}, \quad x * y' \in H$ (où $y'$ est le symétrique de $y$).
	\end{enumerate}
	
	\subsection{Morphismes de groupes}
	Soient $(G, *)$ et $(F, \cdot)$ deux groupes. Un \textbf{morphisme de groupes} de $(G, *)$ dans $(F, \cdot)$ est une application $f: G \to F$ telle que :
	\[ \qcq{x}[, y]{G}, \quad f(x * y) = f(x) \cdot f(y) \]
	Si $f$ est bijective, on dit que c'est un \textbf{isomorphisme} de groupes.
	
	\section{Anneaux et Corps}
	
	\subsection{Anneaux}
	Un \textbf{anneau} est un triplet $(A, +, \cdot)$ où $A$ est un ensemble muni de deux LCI $+$ et $\cdot$ telles que :
	\begin{enumerate}
		\item $(A, +)$ est un groupe commutatif (d'élément neutre noté 0).
		\item La loi $\cdot$ est associative et distributive par rapport à $+$ :
		\[ \qcq{x}[, y][, z]{A}, \quad x \cdot (y + z) = x \cdot y + x \cdot z \quad \text{et} \quad (y + z) \cdot x = y \cdot x + z \cdot x \]
	\end{enumerate}
	Si de plus la loi $\cdot$ admet un élément neutre (noté 1), l'anneau est dit \textbf{unitaire}. Si $\cdot$ est commutative, l'anneau est dit \textbf{commutatif}.
	
	\subsection{Corps}
	Un \textbf{corps} est un triplet $(K, +, \cdot)$ tel que :
	\begin{enumerate}
		\item $(K, +, \cdot)$ est un anneau unitaire commutatif avec $0 \ne 1$.
		\item Tout élément de $K \setminus \{0\}$ est symétrisable pour la loi $\cdot$ (admet un inverse).
	\end{enumerate}
	Exemples : $(\mathbb{Q}, +, \cdot)$, $(\mathbb{R}, +, \cdot)$ et $(\mathbb{C}, +, \cdot)$ sont des corps.
	
	\section{Espaces vectoriels réels}
	
	\subsection{Définition}
	Un \textbf{espace vectoriel réel} est un ensemble $E$ muni d'une loi de composition interne $+$ (loi interne) et d'une loi de composition externe $\cdot$ de $\mathbb{R} \times E$ dans $E$ (loi externe) vérifiant des propriétés de compatibilité :
	\begin{enumerate}
		\item $(E, +)$ est un groupe commutatif.
		\item $\qcq{\lambda}[, \mu]{R}$, $\qcq{\vec{u}}[, \vec{v}]{E}$ :
		\begin{itemize}
			\item $\lambda \cdot (\vec{u} + \vec{v}) = \lambda \cdot \vec{u} + \lambda \cdot \vec{v}$
			\item $(\lambda + \mu) \cdot \vec{u} = \lambda \cdot \vec{u} + \mu \cdot \vec{u}$
			\item $(\lambda\mu) \cdot \vec{u} = \lambda \cdot (\mu \cdot \vec{u})$
			\item $1 \cdot \vec{u} = \vec{u}$
		\end{itemize}
	\end{enumerate}
	Les éléments de $E$ sont appelés des \textbf{vecteurs} et ceux de $\mathbb{R}$ des \textbf{scalaires}.
	
	\subsection{Sous-espaces vectoriels}
	Une partie non vide $F$ de $E$ est un \textbf{sous-espace vectoriel} (sev) de $E$ si $F$ muni des lois induites est un espace vectoriel.
	\textbf{Caractérisation} :
	$F$ est un sous-espace vectoriel de $E$ si et seulement si :
	\begin{enumerate}
		\item $F \ne \emptyset$ (en particulier, $\vec{0}_E \in F$).
		\item $\qcq{\lambda}[, \mu]{R}$, $\qcq{\vec{u}}[, \vec{v}]{F}, \quad \lambda\vec{u} + \mu\vec{v} \in F$.
	\end{enumerate}
	
\end{document}
